\begin{theorem}
\[
\gcd(a^m - 1, a^n - 1) = a^{\gcd(m, n)} - 1
\]
\end{theorem}

\begin{proof}
根据 $\gcd$ 具有 $\gcd(a, b) = \gcd(a - b, b)$ 的性质，不妨设 $m \ge n$，作差有：
\[
\begin{aligned}
\gcd(a^m - 1, a^n - 1) &= \gcd(a^m - a^n, a^n - 1) \\
&= \gcd\bigl(a^n(a^{m - n} - 1), a^n - 1\bigr)
\end{aligned}
\]
由于 $\gcd(a^n, a^n - 1) = 1$，
\[
\begin{aligned}
\gcd\bigl(a^n(a^{m - n} - 1), a^n - 1\bigr)
&= \gcd(a^n, a^n - 1)\,\gcd(a^{m - n} - 1, a^n - 1) \\
&= \gcd(a^{m - n} - 1, a^n - 1)
\end{aligned}
\]
由递归定义，则原式等价于：
\[
\gcd(a^{m \bmod n} - 1, a^n - 1)
\]
根据辗转相除法的形式，我们不难得出：
\[
\gcd(a^m - 1, a^n - 1) = a^{\gcd(m, n)} - 1
\]
\end{proof}

\begin{theorem}
\[
\gcd(a^m - b^m, a^n - b^n) = a^{\gcd(m, n)} - b^{\gcd(m, n)},
\]
其中 $\gcd(a, b) = 1$ 且 $a > b$。
\end{theorem}

\begin{proof}
根据定理1的类似证明方式，作差后有：
\[
\begin{aligned}
&\gcd\bigl(a^m - b^m - (a^n - b^n), a^n - b^n\bigr) \\
&= \gcd\bigl(a^n a^{m-n} - b^n b^{m-n} - (a^n - b^n), a^n - b^n\bigr) \\
&= \gcd\bigl(a^n(a^{m-n} - b^{m-n}) + b^{m-n}(a^n - b^n), a^n - b^n\bigr) \\
&= \gcd\bigl(a^n(a^{m-n} - b^{m-n}), a^n - b^n\bigr) \\
&= \gcd\bigl(a^{m-n} - b^{m-n}, a^n - b^n\bigr) \\
&= \gcd\bigl(a^{m \bmod n} - b^{m \bmod n}, a^n - b^n\bigr)
\end{aligned}
\]
由递归定义可得 $\gcd(a^m - b^m, a^n - b^n) = \bigl| a^{\gcd(m, n)} - b^{\gcd(m, n)} \bigr|$。
\end{proof}

\begin{theorem}
设
\[
G = \gcd\left(\binom{n}{1}, \binom{n}{2}, \ldots, \binom{n}{n-1}\right).
\]
\begin{itemize}
\item 若 $n = p^k$（$k>0$，$p$ 为素数），则 $G = p$；
\item 否则 $G = 1$。
\end{itemize}
\end{theorem}

\begin{proof}
\begin{itemize}
\item 当 $n = p^k$ 时，因为 $\binom{n}{i} = \dfrac{n!}{i!\,(n-i)!}$（$1 \le i < n$）。由勒让德定理，对因子 $p$，$p$ 的出现次数为
\[
L_p(n!) = \sum_{j \ge 1}\left\lfloor\frac{n}{p^j}\right\rfloor.
\]
而
\[
L_p((n-i)!) + L_p(i!) = \sum_{j \ge 1}\left\lfloor\frac{i}{p^j}\right\rfloor + \sum_{j \ge 1}\left\lfloor\frac{n-i}{p^j}\right\rfloor
\le \sum_{j \ge 1}\left\lfloor\frac{n}{p^j}\right\rfloor = L_p(n!).
\]
等号成立当且仅当 $i \mid p^j$ 且 $(n-i) \mid p^j$，但 $n = p^k$ 时不存在这样的 $i$（$1 \le i < n$），因此严格小于。故 $p \mid \binom{n}{i}$ 对所有 $1 \le i < n$ 成立，且 $p \parallel \binom{n}{\lfloor n/2\rfloor}$，从而 $G = p$。

\item 否则，同样根据勒让德定理，有
\[
L_p((n-i)!) + L_p(i!) \le L_p(n!).
\]
设 $n = p_1^{k_1}p_2^{k_2}\cdots p_t^{k_t}$（$t>1$，$p_1<p_2<\cdots<p_t$）。对于每个质因子 $p_0$，将 $n$ 写为 $p_0$ 进制：
\[
n = c_0 + p_0c_1 + p_0^2c_2 + \cdots + p_0^r c_r.
\]
由于 $n$ 不是 $p_0$ 的幂，必存在某个 $\lambda < r$ 使得 $c_\lambda \neq 0$。取 $i = p_0^\lambda c_\lambda$，则可验证等号成立，即
\[
L_{p_0}(n!) = L_{p_0}(i!) + L_{p_0}((n-i)!).
\]
因此对每个素数 $p$ 都存在某个 $i$ 使得 $p \nmid \binom{n}{i}$，从而 $G = 1$。
\end{itemize}
综上，结论成立。
\end{proof}

\begin{theorem}
若 $F(n)$ 表示斐波那契数列第 $n$ 项，则
\[
\gcd(F(n), F(m)) = F(\gcd(n, m)).
\]
\end{theorem}

\begin{proof}
先证明恒等式
\[
F(n+m) = F(n)F(m+1) + F(n-1)F(m).
\]
用归纳法：
\begin{itemize}
\item $m=1$ 时，$F(n+1)=F(n)+F(n-1)$，成立。
\item $m=2$ 时，$F(n+2)=2F(n)+F(n-1)=F(n+1)+F(n)$，成立。
\item 假设对 $m=k-1$ 和 $m=k$ 成立，则对 $m=k+1$ 有
\[
\begin{aligned}
F(n+k-1) &= F(n)F(k) + F(n-1)F(k-1),\\
F(n+k) &= F(n)F(k+1) + F(n-1)F(k).
\end{aligned}
\]
相加得 $F(n+k+1)=F(n)F(k+2)+F(n-1)F(k+1)$。故对任意 $n>1,m>0$ 成立。
\end{itemize}
不妨设 $n \ge m$，则
\[
\begin{aligned}
\gcd(F(n),F(m)) &= \gcd\bigl(F(m+1)F(n-m)+F(m)F(n-m+1),F(m)\bigr)\\
&= \gcd\bigl(F(m+1)F(n-m),F(m)\bigr).
\end{aligned}
\]
由于 $\gcd(F(m+1),F(m)) = \gcd(F(m-1),F(m)) = 1$，故
\[
\gcd(F(m+1)F(n-m),F(m)) = \gcd(F(n-m),F(m)).
\]
由此递归得
\[
\gcd(F(n),F(m)) = \gcd(F(n \bmod m),F(m)) = F(\gcd(n,m)).
\]
\end{proof}

\begin{theorem}
\[
(n+1)\operatorname{lcm}\left(\binom{n}{0}, \binom{n}{1}, \ldots, \binom{n}{n}\right) = \operatorname{lcm}(1,2,\ldots,n+1).
\]
\end{theorem}

\begin{proof}
将 $n+1$ 乘进去：
\[
\operatorname{lcm}\left((n+1)\binom{n}{0}, (n+1)\binom{n}{1}, \ldots, (n+1)\binom{n}{n}\right) = \operatorname{lcm}(1,2,\ldots,n+1).
\]
因此原命题等价于
\[
\operatorname{lcm}\left(n+1, \frac{(n+1)!}{1!(n-1)!}, \ldots, \frac{(n+1)!}{i!(n-i)!}, \ldots, n+1\right) = \operatorname{lcm}(1,2,\ldots,n+1).
\]
左右两边所含最大质因子不超过 $n+1$，故只需对每个质数 $p \le n+1$ 证明两边中 $p$ 的最高幂相等。

由勒让德定理：
\[
L_p((n+1)!) = \sum_{j\ge1}\left\lfloor\frac{n+1}{p^j}\right\rfloor,\qquad
L_p(i!)+L_p((n-i)!) = \sum_{j\ge1}\left\lfloor\frac{i}{p^j}\right\rfloor + \sum_{j\ge1}\left\lfloor\frac{n-i}{p^j}\right\rfloor.
\]
将 $n$ 写为 $p$ 进制：$n = c_0 + pc_1 + p^2c_2 + \cdots + p^r c_r$（$c_r \neq 0$）。欲使 $L_p(i!)+L_p((n-i)!)$ 尽可能小（从而 $\binom{n}{i}$ 中含 $p$ 的幂次最大），类似定理3的讨论，取 $i = p^r c_r - 1$ 分两种情况：
\begin{itemize}
\item 若 $n+1$ 不进位到最高位，则 $p^r \mid i+1$，且
\[
L_p((n+1)!) = L_p((i+1)!)+L_p((n-i)!)=L_p(i!)+L_p((n-i)!)+r.
\]
左边 $p$ 的最高幂为 $p^r$，右边 $n+1$ 的 $p$ 进制表示中 $p^r$ 的系数相同，故相等。
\item 若 $n+1$ 进位，即 $n+1 = p^{r+1}$，取 $i+1 = p^r(p-1)$，$n-i = p^r$。则
\[
L_p((n+1)!) = L_p((i+1)!)+L_p((n-i)!)=L_p(i!)+L_p((n-i)!)+r+1,
\]
左边 $p$ 的最高幂为 $p^{r+1}$，右边 $n+1 = p^{r+1}$，等式成立。
\end{itemize}
因此两边每个素因子的幂次相同，定理得证。
\end{proof}